\begin{textsample}{POS Dim 2 – vem\_ed\_08 – Score 11.00 – vem\_ed\_08/t030.txt}  \label{ex:f2_pos_027}
QUEM É QUEM Vice-diretora- superintendente do Centro Paula Souza (CPS) desde 2019, Emilena Lorenzon Bianco destaca, em \textbf{entrevista} concedida à \textbf{newsletter} VEm, o valor de vivenciar projetos colaborativos com professores de outros países e o que representa o avanço da internacionalização do CPS.
Emilena é doutora em Ciência da Informação (Unesp, 2011), mestra em Engenharia de Produção (UFSCAR, 2004), especialista em Uso Estratégico das Tecnologias da Informação (Unesp, 2001) e graduada em Biblioteconomia (Unesp, 1998).
Suas pesquisas em \textbf{pós-graduação} abordam o cluster calçadista da região de Jaú.
Professora da Fatec Jahu entre 2004 e 2019, atuou na \textbf{implantação} da Fatec Matão (2018) e foi \textbf{coordenadora} de projetos na Inova Paula Souza.
Os Projetos Colaborativos Internacionais (PCIs) são desenvolvidos a distância entre professores e alunos de Fatecs e de instituições de ensino superior (IES) internacionais.
Qual é a importância dos PCIs na formação de profissionais de excelência em âmbito global?
É extremamente importante, tanto para o aspecto profissional, quanto para a vivência pessoal, participar de trabalhos em parceria com outros países.
Graças à tecnologia, que ajudou a \textbf{superar} as barreiras de mobilidade física, atualmente, em todos os \textbf{campos} do conhecimento as práticas são globalizadas e os profissionais atuam de maneira colaborativa, em rede.
Programas como os PCIs oferecem essa oportunidade para nossos professores e alunos.
Por meio da interação virtual, é possível oferecer escalabilidade para experimentação de outras culturas de forma menos custosa, o que permite incluir maior número de alunos.
Essas experiências \textbf{agregam} muito valor ao conhecimento.
Ao ter contato direto com pessoas de outras culturas, os \textbf{jovens} passam a entender a diversidade profissional nos diferentes países e como cada nação tem um jeito de lidar com as \textbf{situações} nos ambientes de trabalho.
Essa riqueza de visões traz inteligência cultural: a gente aprende como \textbf{agir}, como se adaptar e também desenvolve resiliência em \textbf{situações} desafiadoras.
Em trabalho alinhado com a Coordenação de Línguas da CESU, os PCIs contribuem para promover a Internacionalização em Casa.
Qual o papel dos PCIs na internacionalização dos currículos das Fatecs?
Ampliar a representatividade internacional é um objetivo estratégico do Centro Paula Souza.
Para ser reconhecida internacionalmente pela excelência da educação profissional, a instituição precisa formar profissionais

% matched lemmas: agir, agregar, campo, coordenadora, entrevista, implantação, jovem, newsletter, pós-graduação, situação, superar
\end{textsample}
